\chapter{Two full theories at central charge eight}
\label{ch:full_theories}

A chiral algebra organizes left-moving fields. A full local CFT
requires a compatible right-moving theory and a prescription for
pairing their sectors. The Haagerup construction gives two especially
natural choices. They have the same central charge and sharply
different spectra, making them a useful example of why the phrase
``the CFT of a chiral algebra'' needs qualification.

Throughout this chapter we assume the construction inputs and the
structural and positivity conclusions developed in Part II.
The full-theory existence statements use the rational extension
and diagonal constructions. The spectrum calculations themselves
then follow from the character table by elementary counting.

\section{The diagonal theory and a holomorphic extension}
\index{diagonal full theory}\index{holomorphic extension}

The diagonal theory has
\begin{equation}
 \HH_{\rm diag}=\bigoplus_i M_i\otimes\overline{M_i},
 \qquad Z_{\rm diag}=\sum_i|\chi_i|^2.
 \label{eq:full-diagonal}
\end{equation}
The twelve labels $i$ are the simple $U$-modules. They are
self-dual in the present modular category, so charge conjugation
does not change the displayed diagonal pairing. The diagonal
Longo--Rehren construction applied to the completely rational net
of Chapter~\ref{ch:unitarity} gives a unitary local full conformal
net with this Hilbert space. This establishes an actual full theory;
the modular-invariant sum alone would not supply the existence theorem.

The selected Haagerup extension uses the canonical Lagrangian
algebra of $\mathcal H_3$ transported into $\Rep(U)$. Its chiral
state space is
\[
 H'\downarrow_U=\one\oplus b\oplus a^{\oplus2},
 \qquad H'\cong V_{E_8}.
\]
The positive extension is the $E_8$ lattice theory, so its tensor
product with the antiholomorphic theory gives the full CFT
\begin{equation}
 \HH_{E_8}=H'\otimes\overline{H'},\qquad
 Z_{E_8}=|\chi_0+\chi_b+2\chi_a|^2
        =\left|\frac{E_4}{\eta^8}\right|^2.
 \label{eq:full-E8}
\end{equation}
This expression contains off-diagonal $U$-sector pairings. If
$v=(1,1,2,0,\ldots,0)$ in the order $(0,b,a,c_0,\ldots)$, the
coupling matrix is $Z_{ij}=v_iv_j$, whereas the diagonal theory
has $Z_{ij}=\delta_{ij}$.
For example, the diagonal $c_1\otimes\overline{c_1}$ sector supplies
nine scalars of dimension $2/3$, whereas its coefficient in the
extension partition function is zero. Conversely, the extension
contains four copies of $a\otimes\bar a$: each chiral side has two
independent copies of $a$. Changing the pairing changes actual state
multiplicities, not just their names.

Both theories have a unique vacuum and $(c_L,c_R)=(8,8)$.
Holomorphicity of $H'$ means that it has only one simple module
as a chiral algebra. It does not mean that the full theory has
only one state, one local field, or no interesting defects
preserving a smaller algebra. The extension changes which fields
are regarded as chiral-algebra generators.

\section{Reading a full spectrum from chiral grades}
\index{scalar gap}\index{primary multiplicity}

In the diagonal theory, left grade $n$ and right grade $m$ in
sector $i$ give
\begin{equation}
 \begin{split}
 (h,\bar h)&=(h_i+n,h_i+m),\\
 \Delta&=2h_i+n+m,\qquad s_{\rm spin}=n-m,\\
 d_{\rm states}&=a_i(n)a_i(m),\qquad
 d_{\rm primary}=v_i(n)v_i(m).
 \end{split}
 \label{eq:full-grade-count}
\end{equation}
The last line counts fields primary under both Virasoro algebras.
A product of left and right extended-algebra top states is
automatically Virasoro primary; higher grades need descendant
subtraction first.

The scalar condition is $n=m$ in a diagonal sector. Consequently a
positive grade raises a scalar dimension by at least two. Every
relevant scalar, meaning $0<\Delta<2$, therefore comes from a
chiral top with $h_i<1$. Squaring the corresponding top dimensions
gives the complete list:
\begin{center}
\begin{tabular}{c|rrrrrrr}
\toprule
$\Delta$&$2/3$&$10/13$&$12/13$&$14/13$&$16/13$&$4/3$&$22/13$\\
Multiplicity&9&25&49&169&196&729&1764\\
\bottomrule
\end{tabular}
\end{center}
Their total is $2941$. The number looks large relative to twelve
extended sectors because each top space contains multiple local
components. A module with $42$ chiral top states contributes
$42^2=1764$ scalar components when paired diagonally.

\begin{example}[A relevant spinning level]
The $c_1$ module has $h=1/3$, three top states, and
$243-3=240$ new Virasoro primaries at grade one. Pairing a grade-one
primary with a top state gives
\[
 \Delta=\frac13+1+\frac13=\frac53,
 \qquad s_{\rm spin}=\pm1,
\]
with $240\cdot3=720$ components for each sign. These fields are
relevant by dimension but cannot appear alone as rotationally
invariant scalar perturbations.
\end{example}

At dimension two the scalar primaries come from the tops of
$b,a,c_0$ and from current bilinears in the vacuum sector:
\begin{equation}
 d_{\Delta=2,s=0}=80^2+81^2+81^2+6^2=19558.
 \label{eq:full-diagonal-dim2}
\end{equation}
There are also $113$ chiral weight-two primaries at each spin
$\pm2$. Thus the total primary count at dimension two is $19784$,
while the state count is $19798$. The difference is exactly two
stress tensors and twelve current derivatives.

\section{Why the extension has no relevant scalar}

The holomorphic $E_8$ algebra is integer graded. The only state of
weight zero is its vacuum, and the first nontrivial states are the
$248$ currents at weight one. A nonidentity full scalar must have
equal positive left and right weights, so its dimension is at
least two. The first scalar space is
\[
 (H')_1\otimes\overline{(H')_1},
 \qquad\dim=248^2=61504.
\]
Every such current bilinear is a Virasoro primary. Chiral currents
paired with the opposite vacuum have dimension one and spin
$\pm1$, so the ordinary primary gap is one even though the scalar
gap is two.

The two theories can now be compared without ambiguity:
\begin{center}
\begin{tabular}{lrr}
\toprule
Quantity&Diagonal $U$&Selected $E_8$ theory\\
\midrule
Currents per chirality&6&248\\
First nonidentity primary dimension&$2/3$&1\\
First nonidentity scalar dimension&$2/3$&2\\
Relevant scalar components&2941&0\\
Dimension-two scalar components&19558&61504\\
\bottomrule
\end{tabular}
\end{center}
These are exact consequences of the selected characters and full
pairings. They do not identify a lattice Hamiltonian flowing to
either theory.

\section{Which symmetry is being preserved?}
\index{invariant field}

In the diagonal theory, the twelve chiral-algebra-preserving
Verlinde lines form the fusion category underlying
$\Rep(U)=Z(\mathcal H_3)$. The fully invariant bulk space is
$U\otimes\overline U$. Its first nonidentity scalars are the
$6\cdot6=36$ current bilinears of dimension two.

The selected extension has a category of $U$-preserving bimodule
defects equivalent to
$\mathcal H_3\boxtimes\mathcal H_3^{\rm op}$.
Choosing the first chiral factor as the Haagerup symmetry gives
invariant space $U\otimes\overline{H'}$. Its dimension-two scalar
multiplicity is $6\cdot248=1488$. Requiring invariance under both
independent factors gives $U\otimes\overline U$ and multiplicity
$36$. Both choices have invariant scalar gap two, but they are
different symmetry requirements.

This is why a bootstrap comparison must name the symmetry action,
not only the ordinary partition function. Two actions on an
isomorphic full CFT can organize local fields differently. The
defect chapter derives the invariant spaces from explicit line
and tube data.

A dimension-two invariant scalar is only marginal by dimension.
Whether a perturbation is exactly marginal depends on OPE
coefficients, mixing, and higher-order compatibility. Commuting
current constructions provide a standard source of exactly
marginal directions, but the count of all current bilinears is
not the dimension of a conformal manifold.

\begin{example}[A dimension-two scalar that runs]
\label{ex:full-nonabelian-running}
\index{marginal perturbation!running coupling}
Consider a unitary $\mathrm{SU}(2)_k$ current algebra, with
$k$ a positive integer. Fix Cartesian current bases with
\[
 J^a(z)J^b(0)\sim\frac{k\delta^{ab}}{2z^2}
             +\frac{i\varepsilon^{abc}J^c(0)}{z},
 \qquad a,b,c=1,2,3.
\]
Use the same $+i\varepsilon^{abc}$ convention for the independent
antiholomorphic currents. Repeated Lie algebra indices are summed.
The rescaled currents $j^a=\sqrt{2/k}\,J^a$ then satisfy
\[
 j^a(z)j^b(0)\sim\frac{\delta^{ab}}{z^2}
       +i\sqrt{\frac2k}\,
                      \frac{\varepsilon^{abc}j^c(0)}{z}.
\]
The isotropic current bilinear
\[
 \mathcal O=\frac1{\sqrt3}\sum_{a=1}^3j^a\bar j^a
\]
has weights $(1,1)$ and is invariant under simultaneous
$\mathrm{SU}(2)$ rotations of the two current triples.
Its normalization follows directly from the double poles:
\[
 \langle\mathcal O(z,\bar z)\mathcal O(0)\rangle
 =\frac1{3}\sum_{a,b}
           \frac{\delta^{ab}\delta^{ab}}{z^2\bar z^2}
 =\frac1{|z|^4}.
\]

To find its self-OPE coefficient, multiply the two simple poles.
The factor $i^2=-1$ and the identity
$\sum_{a,b}\varepsilon^{abc}\varepsilon^{abd}
=2\delta^{cd}$ give
\[
 -\frac{2}{3k|z|^2}
       \sum_{a,b,c,d}\varepsilon^{abc}\varepsilon^{abd}
                  j^c\bar j^d
 =-\frac{4}{3k|z|^2}\sum_cj^c\bar j^c
 =-\frac{4}{\sqrt3\,k}\frac{\mathcal O}{|z|^2}.
\]
Thus
\begin{equation}
 \mathcal O(z,\bar z)\mathcal O(0)
 =\frac{\one}{|z|^4}
  +C_{\mathcal O\mathcal O\mathcal O}
                          \frac{\mathcal O(0)}{|z|^2}
  +\cdots,
 \qquad
 C_{\mathcal O\mathcal O\mathcal O}
       =-\frac4{\sqrt3\,k}.
 \label{eq:full-isotropic-OPE}
\end{equation}
Terms with one double pole and one simple pole vanish because
$\delta^{ab}\varepsilon^{abc}=0$. The omitted terms include
spinning fields and higher-dimension fields. In this invariant
current sector, the only dimension-two scalar generated by the
displayed collision is $\mathcal O$ itself.

Now fix the action and scale conventions, since these determine
the sign of the beta function. Use the Euclidean weight
$e^{-S_g}$ with
\[
 S_g=S_0+g\int d^2x\,\mathcal O(x),\qquad
 d^2x=dx_1\,dx_2.
\]
The velocity has been absorbed into the Euclidean time coordinate.
There is no extra $1/(2\pi)$ in the definition of $g$.
Let $a$ be the short-distance cutoff and increase the coarse-graining
length by $e^{d\ell}$. The second-order term in $e^{-S_g}$ is
proportional to $+g^2/2$. Its short-distance shell contribution is
\begin{align*}
 \frac{g^2}{2}\int d^2y
       \int_{a<|r|<ae^{d\ell}}d^2r\,
       \frac{C_{\mathcal O\mathcal O\mathcal O}}{|r|^2}
                         \mathcal O(y)
 &=
 \pi C_{\mathcal O\mathcal O\mathcal O}g^2d\ell
                         \int d^2y\,\mathcal O(y),\\
 \int_{a<|r|<ae^{d\ell}}\frac{d^2r}{|r|^2}
 &=2\pi\,d\ell.
\end{align*}
The factor $1/2$ prevents double counting the two integrated
insertions. Matching the term linear in $\mathcal O$ to
$-\bigl(g+dg\bigr)\int\mathcal O$ therefore gives
\begin{equation}
 \beta_{\rm length}(g):=\frac{dg}{d\ell}
 =-\pi C_{\mathcal O\mathcal O\mathcal O}g^2+O(g^3)
 =B_k g^2+O(g^3),
 \qquad B_k=\frac{4\pi}{\sqrt3\,k}>0.
 \label{eq:full-isotropic-beta}
\end{equation}
There is no linear term because $\Delta_{\mathcal O}=2$.
The identity collision shifts the vacuum energy, while a circular
shell averages the nonzero-spin terms to zero. A beta function
defined using an inverse length $\mu$ has the opposite sign:
$\mu\,dg/d\mu=-\beta_{\rm length}(g)$.

Dropping the cubic and higher terms gives the leading running,
\[
 g(\ell)=\frac{g_0}{1-B_k g_0\ell},
 \qquad \ell=\log(r/a).
\]
For $g_0<0$, the coupling approaches zero from below and is
\emph{marginally irrelevant}. In the asymptotic regime,
\[
 g(r)=-\frac1{B_k\log(r/a)}
       +O\!\left(\frac{\log\log(r/a)}{\log^2(r/a)}\right),
\]
where the first subleading terms also encode the initial coupling
and higher-order beta function. A dimensionless observable with
an expansion $Q(g)=Q_*+q_1g+O(g^2)$ consequently approaches its
limit as $1/\log r$ when $q_1\ne0$. If that linear coefficient
vanishes, its leading drift can be smaller; one should not fit
every observable with the same logarithm.
For $g_0>0$, the coupling grows and is marginally relevant.
The leading solution indicates loss of weak-coupling control at
a length of order $a\exp[1/(B_k g_0)]$. Its pole does not determine
a physical infrared endpoint.

Compare the unit-normalized Cartan scalar
$\mathcal O_{\rm C}=j^3\bar j^3$. Since
$\varepsilon^{33c}=0$, its self-OPE has no
$\mathcal O_{\rm C}/|z|^2$ term, so this quadratic beta-function
obstruction is absent. Exact marginality requires more than that
one vanishing coefficient: the abelian current-current theorem,
under its unitary commuting-current hypotheses, supplies the
all-order local conformal family. The isotropic sum uses
noncommuting currents and has the nonzero obstruction just
computed. Both perturbing fields have dimension two at the
starting CFT, but their renormalization-group behavior differs.
\end{example}

\begin{example}[The same running produces two different logarithms]
\label{ex:full-current-logarithms}
\index{logarithmic corrections}
For a concrete comparison model, take the standalone diagonal
$\mathrm{SU}(2)_1$ WZW theory and perturb it by the preceding
isotropic scalar with $g_0<0$. Its spin-$(1/2,1/2)$ primary has
$(h,\bar h)=(1/4,1/4)$. Under the diagonal global
$\mathrm{SU}(2)$ its four components split into a singlet and
a triplet, labelled by $J=0,1$. Both have unperturbed scaling
dimension $\Delta_0=1/2$.
These are internal $\mathrm{SU}(2)$ representation spins;
all four components have conformal spin $h-\bar h=0$.

The current Ward identities turn the simple poles with a primary
into the left and right spin generators $t_L^a,t_R^a$.
Before specializing the level, the coefficient of that primary
in $\mathcal O(z,\bar z)\Phi_J(0)$ is
\[
 \mathcal O(z,\bar z)\Phi_J(0)
 \sim\frac{q_J}{|z|^2}\Phi_J(0),\qquad
 q_J=\frac{2}{\sqrt3\,k}\,t_L\cdot t_R
     =\frac{J(J+1)-3/2}{\sqrt3\,k}.
\]
The last equality is ordinary angular-momentum addition:
$2t_L\cdot t_R=(t_L+t_R)^2-t_L^2-t_R^2$, with
$t_L^2=t_R^2=3/4$. The singlet and triplet diagonalize the
first-order mixing.

With the action convention above, first-order perturbation theory
for $\langle\Phi_J(r)\Phi_J^\dagger(0)\rangle$ inserts
$-g\int\mathcal O$. The two small neighborhoods of the external
fields each contribute $2\pi q_J\log(r/a)$. Thus the correlator
has relative correction $-4\pi gq_J\log(r/a)$.
Comparing with
$(r/a)^{-2(\Delta_0+\delta\Delta_J)}
=(r/a)^{-2\Delta_0}
 \bigl[1-2\delta\Delta_J\log(r/a)+\cdots\bigr]$
gives
\[
 \delta\Delta_J=2\pi gq_J+O(g^2).
\]
Replacing $g$ by its running value resums the leading logarithms:
\begin{align*}
 G_J(r)&\propto r^{-2\Delta_0}
  \exp\!\left[-4\pi q_J
         \int_0^{\log(r/a)}g(\ell)\,d\ell\right]\\
 &\propto r^{-2\Delta_0}
     \bigl[1+B_k|g_0|\log(r/a)\bigr]^{4\pi q_J/B_k}.
\end{align*}
The amplitudes and the reference scale inside the logarithm
depend on the normalization and initial coupling. The leading
powers follow from the two coefficients already calculated:
\[
 \frac{4\pi q_J}{B_k}=J(J+1)-\frac32.
\]
At $k=1$ this gives, to leading logarithmic accuracy,
\[
 G_0(r)\propto\frac1{r[\log(r/a)]^{3/2}},
 \qquad
 G_1(r)\propto\frac{[\log(r/a)]^{1/2}}{r}.
\]
Two fields with the same limiting dimension can therefore drift
in opposite directions and carry different multiplicative
logarithms. This is a canonical WZW comparison calculation.
Applying it to a particular lattice chain would require identifying
its continuum currents, perturbation, and coupling sign. The
calculation neither identifies the Haagerup chain nor fixes a
gapped endpoint, and it does not transport Haagerup defect
junctions along a deformation.
\end{example}

\section{Four explicit marginal directions and a radius example}
\index{current-current deformation}\index{Zamolodchikov metric}\index{Narain moduli}

We can say more than the general warning. The two Cartan currents
in $U$ commute and can be normalized as
\[
 j^1=\frac{H_L}{\sqrt2},\qquad
 j^2=\frac{H_R}{\sqrt{78}},\qquad
 j^a(z)j^b(0)\sim\frac{\delta^{ab}}{z^2}.
\]
The subscripts on $H_L,H_R$ still label the two internal
affine factors, both holomorphic. Introduce an independent
antiholomorphic copy $\bar j^b$. The four fields
\begin{equation}
 \mathcal O_{ab}=j^a\bar j^b,\qquad a,b=1,2,
 \label{eq:full-abelian-marginals}
\end{equation}
are scalar primaries of dimension two, with
\[
 \langle\mathcal O_{ab}(z,\bar z)\mathcal O_{cd}(0)\rangle
 =\frac{\delta^{ac}\delta^{bd}}{|z|^4}.
\]
Thus their seed-point Zamolodchikov metric, the two-point metric
on infinitesimal deformations in this normalization, is the
identity matrix.

The abelian current-current deformation theorem supplies a
simultaneous local exactly marginal family in these four
directions. Its hypotheses involve mutually commuting left
currents and mutually commuting right currents with the
positive current two-point forms. The vanishing structure
constants explain the absence of the usual quadratic
current-current beta-function obstruction; the theorem controls
the all-order statement. One should not infer that statement
from the vanishing of a single three-point coefficient alone.

There is a useful geometric description. For $r$ commuting
currents on each side, collect a state's abelian charges in
$p=(p_L,p_R)\in\R^{r,r}$. The charge zero-mode contributions
are $p_L^2/2$ and $p_R^2/2$. A general state also has
Heisenberg oscillator levels and neutral-sector weights; these
pieces stay fixed in the charge-boost description.
An infinitesimal boost
\begin{equation}
 \delta p_L=\varepsilon p_R,\qquad
 \delta p_R=\varepsilon^{\mathsf t}p_L,
 \label{eq:full-charge-boost}
\end{equation}
preserves $p_L^2-p_R^2$ and therefore the spin contribution.
It changes the total dimension by
$\delta\Delta=2p_L^{\mathsf t}\varepsilon p_R$.
Independent rotations of left and right current bases do not
change the theory's charge geometry. The associated local
deformation space is consequently
\[
 O(r,r)/(O(r)\times O(r)),
 \qquad\dim=r^2.
\]
To see the count, an infinitesimal $O(r,r)$ transformation has blocks
\[
 X=\begin{pmatrix}A&B\\B^{\mathsf t}&D\end{pmatrix},
 \qquad A^{\mathsf t}=-A,\quad D^{\mathsf t}=-D.
\]
The antisymmetric blocks rotate the two current bases. After removing
them, the arbitrary real $r$-by-$r$ matrix $B$ supplies $r^2$ directions.
For the chosen two Cartans this gives four local directions.
This description does not assert that $U$ is a tensor product
of two free bosons with a neutral algebra; the ambient theory
can impose nontrivial charge-sector selection rules.
Nor does it determine the global duality quotient.

\begin{example}[An enhanced current leaves the chiral algebra]
In the compact-boson convention of Chapter~\ref{ch:cft},
let $R=\sqrt2e^t$. The state with $(m,n)=(1,1)$ has
\[
 p_L=\sqrt2\cosh t,\qquad p_R=-\sqrt2\sinh t,
\]
so
\[
 h=\cosh^2t,\qquad \bar h=\sinh^2t,\qquad
 s_{\rm spin}=1,\qquad\Delta=\cosh(2t).
\]
At $t=0$ it is a holomorphic weight-one current. For $t\ne0$
it is a nonchiral local field of the same integer spin, while
the theory remains conformal with $c=1$. This is how an exactly
marginal deformation can remove an enhanced chiral symmetry
without destroying conformal invariance. In charge space,
the radius change is the boost
$p_L(t)=\cosh t\,p_L(0)-\sinh t\,p_R(0)$ and its
right-moving counterpart.
\end{example}

For the ordinary full $E_{8,1}$ theory, the eight Cartans on each
side give the familiar $64$ toroidal directions, locally
$O(8,8)/(O(8)\times O(8))$. The special rational point has
many more dimension-two current bilinears than these $64$
unobstructed directions. The same distinction applies to the
$36$ bilinears in the diagonal $U$ theory.

The four fields in \eqref{eq:full-abelian-marginals} are
categorical singlets at the rational point. Continuing a full
system of topological lines along a deformation also requires
transporting their junctions and checking locality with the
deformed fields. Exact marginality by itself does not prove
that the specified Haagerup category survives everywhere on
the family. This is a question about the deformed symmetry
data, separate from the existence of the ordinary conformal
deformation.

\section{Thermodynamics from the same spectrum}
\index{thermal trace}\index{thermal tail bound}

Let $t=\beta v/L$ and $b=2\pi t$. For either full theory,
define the dimensionless energy
$E=\Delta-2/3$, since $(c_L+c_R)/24=2/3$. Then
\begin{equation}
 Z(b)=\sum_{\Delta,s}d_{\Delta,s}e^{-bE},\qquad
 H=\frac{2\pi v}{L}E.
 \label{eq:full-thermal}
\end{equation}
The vacuum energy is $-4\pi v/(3L)$. Once $Z$ converges,
\[
 \langle E\rangle=-\partial_b\log Z,\qquad
 C_V=b^2\bigl(\langle E^2\rangle-\langle E\rangle^2\bigr),
 \qquad S=\log Z+b\langle E\rangle.
\]
A finite spectral truncation gives a lower bound on $Z$ because
all degeneracies are nonnegative. To make it a quantitative
approximation, one needs an upper bound on the omitted tail.

\begin{proposition}[A modular tail bound]
\label{prop:full-thermal-tail}
Assume a convergent positive modular-invariant partition function
$Z(b)=Z(4\pi^2/b)$. Suppose every state with $\Delta\le D$
is known, where $D>2/3$. Write $Z=L_D+H_D$ for the known and
omitted sums, and put $E_D=D-2/3$. If $b>2\pi$ and
$b'=4\pi^2/b$, then
\begin{equation}
 0\le H_D(b)\le
 \frac{L_D(b)-L_D(b')}{e^{(b-b')E_D}-1}.
 \label{eq:full-tail}
\end{equation}
At $b=2\pi$, provided the differentiated trace converges,
$0\le H_D(b)\le L_D'(b)/E_D$.
\end{proposition}
\begin{proof}
Every omitted state has $E>E_D>0$, so
$H_D(b)\le rH_D(b')$ with
$r=e^{-(b-b')E_D}<1$. Modular invariance gives
$H_D(b')=H_D(b)+L_D(b)-L_D(b')$. Substitution and division by
$1-r$ prove the first claim.
At the fixed point, differentiating the modular identity gives
$Z'(2\pi)=0$. Hence
$L_D'=-H_D'\ge E_DH_D$.
\end{proof}

The estimate is an example of a global consistency relation
controlling information beyond a finite table. It requires
modular invariance and positivity as hypotheses; finite agreement
of a few coefficients cannot establish either one.
For moment tails choose $0<b_0<b$ and use
\begin{equation}
 \sum_{\Delta>D,s}d_{\Delta,s}E^ke^{-bE}
 \le H_D(b_0)\sup_{E\ge E_D}E^ke^{-(b-b_0)E}.
 \label{eq:full-moment-tail}
\end{equation}
For $k>0$, the supremum is at $E_D$ if
$(b-b_0)E_D\ge k$, otherwise at $k/(b-b_0)$.
Thus one can control the energy and heat capacity as well as
the partition function.

\begin{example}[A universal value at the modular fixed point]
At $b=2\pi$, modular invariance gives
$\langle E\rangle=0$. Therefore
\[
 \langle\Delta\rangle=\frac23,
 \qquad S=\log Z
\]
at $t=1$ in either theory. The same mean dimension does not imply
the same distribution, entropy, or heat capacity. The two spectra
give different values of these other observables.
\end{example}

\section{What has and has not been reconstructed}

We now know exact state counts, Virasoro-primary multiplicities,
several symmetry-resolved gaps, and a way to obtain controlled
thermal observables. These are much more than a central charge.
They remain less than a complete list of local OPE constants.
The next chapter computes actual chiral correlators in the lattice
realization, including a normalized four-point function that can
be expanded by hand. The sewing chapter will explain the additional
data needed for nonvacuum bulk correlators of the diagonal theory.

\section{Exercises}

\begin{exercise}\label{ex:full-relevant}
Reproduce the seven relevant scalar multiplicities and their sum
$2941$. Then compute the spinning multiplicities at dimensions
$23/13$ and $25/13$ using the $d_4$ and $d_1$ rows.

\par\smallskip\noindent{\footnotesize\normalfont
\hyperref[sol:ex:full-relevant]{Solution on page~\pageref*{sol:ex:full-relevant}.}}\par
\end{exercise}

\begin{exercise}\label{ex:full-E8states}
Show that the full $E_8$ theory has $69752$ states at dimension
two. Separate them into scalar primaries, chiral primaries,
current derivatives, and stresses.

\par\smallskip\noindent{\footnotesize\normalfont
\hyperref[sol:ex:full-E8states]{Solution on page~\pageref*{sol:ex:full-E8states}.}}\par
\end{exercise}

\begin{exercise}\label{ex:full-tail}
Derive \eqref{eq:full-moment-tail} and evaluate its supremum for
$k=1,2$. Explain why a reliable heat-capacity error bound needs
more than an error bound on $Z$ alone.

\par\smallskip\noindent{\footnotesize\normalfont
\hyperref[sol:ex:full-tail]{Solution on page~\pageref*{sol:ex:full-tail}.}}\par
\end{exercise}

\begin{exercise}\label{ex:full-onevac}
Suppose a proposed branching prescription gives
$Z=6|\chi_0|^2+\cdots$ with all coefficients nonnegative and
all omitted sectors of positive weight. Explain the vacuum
problem. Why does dividing the expression by six not automatically
produce a new local CFT?

\par\smallskip\noindent{\footnotesize\normalfont
\hyperref[sol:ex:full-onevac]{Solution on page~\pageref*{sol:ex:full-onevac}.}}\par
\end{exercise}

\begin{exercise}\label{ex:full-symmetry}
Explain the three different meanings of an invariant scalar count
$36$, $1488$, or $61504$ in this chapter. Which of these numbers
uses no symmetry restriction at all?

\par\smallskip\noindent{\footnotesize\normalfont
\hyperref[sol:ex:full-symmetry]{Solution on page~\pageref*{sol:ex:full-symmetry}.}}\par
\end{exercise}

\begin{exercise}\label{ex:full-boost}
Let $\eta=\diag(I_r,-I_r)$ and
$X=\left(\begin{smallmatrix}0&\varepsilon\\
\varepsilon^{\mathsf t}&0\end{smallmatrix}\right)$.
Verify $X^{\mathsf t}\eta+\eta X=0$ and derive
\eqref{eq:full-charge-boost}. For $r=1$, recover the radius
example and show explicitly why a chiral current can acquire
$\bar h>0$ without losing integer spin.

\par\smallskip\noindent{\footnotesize\normalfont
\hyperref[sol:ex:full-boost]{Solution on page~\pageref*{sol:ex:full-boost}.}}\par
\end{exercise}
